\section{引言}
远程状态估计通过网络利用传感器测量形成系统状态预测。
随机事件触发机制根据当前量测与预测的偏差决定是否发送，静默事件也参与条件估计。
指数二次触发规则可在适当线性Gaussian模型下保持条件Gaussian结构，为通信和估计分析提供基础\cite{han2015}。
当发送载荷遭到修改时，远端估计偏差还可能通过实际预测反馈改变未来发送概率。
本文研究这一可信事件路径能够对真实估计损伤提供怎样的约束。

已有安全估计研究已经分析统计隐蔽与性能退化的关系\cite{bai2017,shang2022}。
Shang等研究随机事件估计下的创新分布与最坏攻击\cite{shang2022}；
Guo等在明确信息、调度和条件矩约束下比较包括非线性策略在内的可行策略\cite{guo2025}。
Song和Ye通过更改预测反馈调节随机发送概率，并优化指定隐蔽条件下的性能退化\cite{song2022}。
Xu等和Wei等分别利用可靠与可疑信息源、历史创新及多级检测约束研究最坏损伤\cite{xu2023,wei2026}。
Wei等还在可靠数据不可读取且相应交叉矩满秩时，给出不存在非平凡可行隐蔽攻击的条件。
因此，随机触发、可信信息、历史数据与隐蔽损伤本身均已有研究基础。

本文考虑另一项明确的信任条件：发送事件与远端实际预测反馈可信，修改仅发生在真实发送决定之后的载荷上。
反馈忠实于远端实际算出的值，但该值可能已经包含先前修改造成的偏差。
审计者没有源端真实测量的独立副本，只能使用预先规定的记录。
在这一条件下，需要回答：记录与正常运行在分布上接近，是否足以限制真实MSE损伤；若足够，需要观察到哪个时刻、保留哪些历史校准信息？
TV提供对任意同记录检验的分布区分解释，但TV本身不能控制无界平方损失。
所需的额外关系来自隐藏估计偏差对未来可信事件概率的影响。

数字反馈进一步引出精度问题。前向量化观测下的篡改、反馈辅助的创新修改和调度防护已有直接成果\cite{li2026,li2025}，
随机事件与量化创新的联合估计也已有研究\cite{shen2025}。
本文扩展仅量化实际预测反馈中心，保持发送载荷连续，并将估计器内部信息与外部审计接口分开。
量化的多对一映射可使内部预测改变而下一反馈码不变，因而削弱指定记录所提供的保证。
这不等于任何量化估计器都存在相同安全限制。

本文主线为一项精确协议的损伤保证及其精度边界。
首先，建立历史校准的下一静默统计量与隐藏偏差之间的统一关系，并保留事件选择造成的协方差期望变化；
由此给出有限时域双预算最坏损伤的匹配阶数。
其次，作为正文扩展，在标量量化反馈模型中证明受限可信记录同律时的残余损伤及其幅值--精度尺度。
该扩展的统一概率下界适用于每个固定有限目标时刻，其常数不要求对时间一致。
单拍符号构造、Gaussian积分和TV不等式均作为技术工具，不单列为方法创新。

本文采用受限可信遥测作为明示的信息结构，不声称已验证某工业产品满足全部假设。
实际系统可以配置仅发布指定字段的接口\cite{opcdata}，但端点信任、原始数据访问和验证失败处理仍需按部署确定。
本文不证明部分保护比完整载荷认证更便宜，也不将风险分析替代完整性防护。
